Penelitian direncanakan berlangsung selama enam bulan. Rincian kegiatan dan alokasi waktunya disajikan pada Tabel~\ref{tbl:timeline}.

\newcommand{\G}{\cellcolor{green}}   % 50% green, lighter
\newcommand{\B}{\cellcolor{blue!70}}    % light blue
\newcommand{\Y}{\cellcolor{yellow}}  % softer yellow
\begin{table}[H]
  \centering
  \caption{Jadwal Pelaksanaan}
  \label{tbl:timeline}
  \renewcommand{\arraystretch}{1.3}
  \resizebox{\textwidth}{!}{%
  \begin{tabular}{|c|p{6cm}|c|c|c|c|c|c|}
  \hline
  \textbf{No.} & \textbf{Jenis Kegiatan} & \textbf{Bulan 1} & \textbf{Bulan 2} & \textbf{Bulan 3} & \textbf{Bulan 4} & \textbf{Bulan 5} & \textbf{Bulan 6} \\
  \hline
  1 & Identifikasi masalah & \B &  &  &  &  &  \\
  \hline
  2 & Studi literatur & \B & \B &  &  &  &  \\
  \hline
  3 & Perancangan model sistem ANC dan skenario uji &  & \B &  &  &  &  \\
  \hline
  4 & Implementasi algoritma FXLMS dan FXNLMS pada simulasi &  & \B & \G &  &  &  \\
  \hline
  5 & Simulasi dan evaluasi kinerja peredaman &  &  & \G &  &  &  \\
  \hline
  6 & Analisis pengaruh normalisasi \emph{step-size} &  &  & \G & \G &  &  \\
  \hline
  7 & Perancangan dan implementasi algoritma pada FPGA &  &  &  & \G & \Y &  \\
  \hline
  8 & Sintesis dan evaluasi utilisasi sumber daya &  &  &  &  & \Y &  \\
  \hline
  9 & Pengujian prototipe dan pengukuran peredaman &  &  &  &  & \Y & \Y \\
  \hline
  10 & Perbandingan efisiensi komputasi FXLMS dan FXNLMS &  &  &  &  & \Y & \Y \\
  \hline
  11 & Analisis akhir &  &  &  &  & \Y & \Y \\
  \hline
  12 & Pembuatan laporan &  &  &  &  & \Y & \Y \\
  \hline
  \end{tabular}
  } % akhir resizebox
\end{table}
